% =====================================================================
%  Appendices. Included by factor-model-paper.tex.
% =====================================================================
\appendix
\newpage

\section{Statistical methodology}
\label{app:method}

This appendix states each procedure in the form in which it is implemented, including
the conventions --- bandwidth rules, degrees of freedom, boundary handling --- that
decide whether two nominally identical implementations agree. Every routine named
here is a pure function over arrays in \code{backend/core} and is exercised by the
test suite against a process whose answer is known analytically.

\subsection{Conventions}

The significance level is $\alpha = 0.05$ throughout, and every $p$-value is stored
and displayed so that a reader may disagree with the threshold. Returns are natural
logarithms. Annualisation uses 252 trading days: a daily standard deviation is
multiplied by $\sqrt{252}$ and a daily mean by 252. Sample standard deviations use
$\mathrm{ddof}=1$ unless stated; the Ledoit-Wolf derivation is the exception and is
noted where it arises.

% ---------------------------------------------------------------------
\subsection{The stationarity battery}
\label{app:stationarity}

A single Dickey-Fuller test is useless as a gate on daily returns: it rejects almost
always, so it cannot discriminate. The battery is therefore built around what
actually goes wrong with financial series --- an untransformed level, a structural
break, stale pricing, illiquidity --- with each test given the role its own null
hypothesis supports.

\subsubsection{Unit-root tests, read jointly}

\begin{description}[leftmargin=1.6em,style=nextline,itemsep=0.35ex]
\item[Augmented Dickey-Fuller \citep{dickey1979}.] $H_0$: a unit root is present. Lag
  order is selected by AIC; the regression includes a constant, which is the right
  specification for an excess return with a possible non-zero mean.
\item[KPSS \citep{kpss1992}.] $H_0$: the series \emph{is} stationary --- the exact
  complement of ADF. Level stationarity, with the automatic Newey-West bandwidth.
\item[Phillips-Perron \citep{phillips1988}.] Same null as ADF but with a
  non-parametric correction for serial correlation and heteroskedasticity in place of
  ADF's augmentation lags. Reported alongside as a robustness check, not gated on.
\end{description}

ADF and KPSS have opposite nulls, and reading them together yields four verdicts where
either alone yields two. This is the single most useful idea in the battery:

\begin{table}[htbp]
\centering
\small
\caption{Joint reading of ADF and KPSS at $\alpha = 0.05$.}
\begin{tabular}{@{}llll@{}}
\toprule
\textbf{ADF} & \textbf{KPSS} & \textbf{Verdict} & \textbf{Consequence} \\
\midrule
rejects & does not reject & stationary & pass \\
does not reject & rejects & \textbf{unit root} & \textbf{fail} --- must be differenced \\
rejects & rejects & break or heteroskedasticity & warn \\
does not reject & does not reject & inconclusive --- usually too short & warn \\
\bottomrule
\end{tabular}
\end{table}

Only the second row blocks a series. The third is the interesting one: both tests
rejecting is not a contradiction but a signature --- it indicates a structural break or
strong heteroskedasticity rather than a clean unit root, and it is the only condition
under which the Zivot-Andrews test is run.

\subsubsection{Structural break}

\textbf{Zivot-Andrews \citep{zivot1992}} tests a unit root against stationarity around
a single endogenously determined break, searching every candidate break date and
taking the minimum statistic. It is by far the most expensive test in the battery,
which is why it runs only when the joint reading says ``not a clean unit root'' and at
least 100 observations are available. The estimated break date is stored, because a
break in 2008-09 and a break in 2020-03 invite different conclusions.

\subsubsection{Stale pricing}

\textbf{Lo-MacKinlay variance ratio \citep{lo1988}}, heteroskedasticity-robust, at
lags $q \in \{2,5,10\}$:

\begin{equation}
\mathrm{VR}(q) \;=\; \frac{\operatorname{Var}\!\left(\sum_{j=0}^{q-1} r_{t-j}\right)}{q \operatorname{Var}(r_t)} .
\end{equation}

Under a random walk $\mathrm{VR}=1$. Values above 1 indicate positive autocorrelation
--- trending or smoothed pricing --- and values below 1 mean reversion. The
heteroskedasticity-robust variant is essential: daily returns are conditionally
heteroskedastic by construction, and the homoskedastic version rejects on that alone.

\begin{decision}{A level, not a return}
The \code{arch} package's variance-ratio implementation consumes a \emph{price level}
and differences internally. Passing returns to it differences twice, giving
$\mathrm{VR} \approx 1/q$ and flagging every clean factor as mean-reverting. The
implementation here therefore reconstructs a level by cumulative summation before
calling it, and a test asserts that a smoothed series shows $\mathrm{VR} > 1$ --- which
is how the defect was found.
\end{decision}

\textbf{Ljung-Box \citep{ljung1978}} at 10 lags, $H_0$: no autocorrelation up to lag
10. It is the second stale-pricing signal, and is also computed on squared returns,
where rejection indicates volatility clustering rather than predictability in the mean.

\subsubsection{Conditional heteroskedasticity}

\textbf{Engle's ARCH-LM \citep{engle1982}} at 10 lags, $H_0$: no conditional
heteroskedasticity. \emph{This test is recorded and never gated on.} A GARCH process
is strictly stationary; gating on ARCH-LM would exclude essentially every financial
series ever published. The reason to know about clustering is that it argues for HAC
standard errors --- not against the factor.

\subsubsection{Distribution and liquidity}

Skewness, excess kurtosis (Fisher convention, zero for a Gaussian) and the
Jarque-Bera statistic are recorded. Like ARCH-LM, Jarque-Bera is never gated on: fat
tails are the normal condition of a daily return series.

The \textbf{zero-return share} is the cheapest and frequently the most telling check in
the battery. Above 50\% the series is not really trading daily and the verdict fails;
above 20\% it warns. The \textbf{maximum absolute return} is recorded as a
data-integrity check --- a single print of $-5.775$ in log space is not a market event
but a corporate action or a bad first row.

\subsubsection{Verdict assembly}

\begin{center}
\small
\begin{tabular}{@{}lll@{}}
\toprule
\textbf{Condition} & \textbf{Verdict} & \textbf{Flag} \\
\midrule
Joint reading is ``unit root'' & fail & \code{unit\_root} \\
Fewer than 60 observations & fail & \code{too\_short} \\
Zero variance & fail & \code{constant} \\
Zero-return share $> 0.50$ & fail & \code{illiquid} \\
\addlinespace
Joint reading inconclusive & warn & \code{inconclusive\_unit\_root} \\
Joint reading is ``break'' & warn & \code{structural\_break} \\
Fewer than 252 observations & warn & \code{short\_history} \\
Zero-return share $> 0.20$ & warn & \code{sparse\_trading} \\
$|\mathrm{VR}-1| > 0.25$ with a significant $p$ & warn & \code{stale\_pricing} \\
\bottomrule
\end{tabular}
\end{center}

Every flag carries a human-readable reason, and the reason rather than the flag is
what the interface displays.

\subsubsection{Sparse release factors}

A release-event factor is zero on roughly 80\% of days by construction. The liquidity
and stale-pricing gates would read that as a dead instrument. For a series marked
sparse, the tests therefore run on the non-zero subsequence --- the statistically
meaningful object is the sequence of releases --- and the sparsity is reported rather
than penalised. A short window over such a factor holds few releases, so it warns
instead of failing, with the reason directing the reader to the full-history
diagnostic where power is adequate.

% ---------------------------------------------------------------------
\subsection{Regression}
\label{app:regression}

\subsubsection{Weighted least squares}

With design matrix $X$ (a column of ones followed by the factor panel) and weights $w$
normalised to mean one, the estimator solves the weighted normal equations. Equal
weighting sets $w_t = 1$; exponential weighting sets

\begin{equation}
w_t \;\propto\; 2^{-(T-t)/H},
\end{equation}

for half-life $H$ in days, again normalised to mean one so that the effective sample
size is comparable across settings.

\subsubsection{Huber}

Robust regression \citep{huber1964} by iteratively reweighted least squares, five
iterations, with tuning constant $\delta = 1.345$ --- the value giving 95\% asymptotic
efficiency under Gaussian errors. The weight function is

\begin{equation}
\psi(z) \;=\; \begin{cases} 1 & |z| \le \delta \\ \delta/|z| & |z| > \delta\end{cases},
\qquad z_t = \frac{\varepsilon_t}{\hat s},
\end{equation}

with $\hat s$ a robust scale estimate of the residuals. Observations inside $\delta$
robust standard deviations keep unit weight; those outside are downweighted in
proportion to their distance rather than discarded.

\subsubsection{Ridge}
\label{app:ridge}

The penalty applies to standardised slopes so that $\lambda$ does not depend on
whether a factor is quoted in percent or in decimals, and the intercept is left
unpenalised so that $\alpha$ is not biased toward zero. Penalising
$\sum_j (b_j \sigma_j)^2$ gives normal equations

\begin{equation}
\left(X^{\prime}X + \lambda n \operatorname{diag}(\sigma^2)\right) b \;=\; X^{\prime}y ,
\end{equation}

so the penalty is \emph{multiplied} by $\sigma^2$.

\begin{decision}{The sign of an exponent}
The original implementation divided by $\sigma^2$ instead. For daily returns with
$\sigma \approx 0.01$ that inflates the penalty by a factor of $10^4$, and $\lambda=1$
shrank every loading to exactly zero. The regression still returned, still reported an
$R^2$, and still produced a loading vector; it was simply the zero vector. A test that
asserts a modest ridge penalty leaves a known loading recognisable is what caught it.
\end{decision}

\subsubsection{HAC standard errors}

Newey-West \citep{newey1987} with Bartlett (triangular) weights:

\begin{equation}
V \;=\; (X^{\prime}X)^{-1}\left[S_0 + \sum_{\ell=1}^{L} w_\ell\left(\Gamma_\ell + \Gamma_\ell^{\prime}\right)\right](X^{\prime}X)^{-1},
\qquad w_\ell = 1 - \frac{\ell}{L+1},
\end{equation}

where $\Gamma_\ell$ is the lag-$\ell$ autocovariance of the score vectors
$\varepsilon_t x_t$. Bartlett weights guarantee a positive semi-definite result, which
a truncated (unweighted) kernel does not. The automatic bandwidth is

\begin{equation}
L \;=\; \left\lfloor 4\left(\frac{T}{100}\right)^{2/9} \right\rfloor ,
\end{equation}

giving 4 lags at 100 observations, 5 at 252 and 6 at 1000. The bandwidth is
deliberately small: autocorrelation in daily returns is short-lived, and the cost of
over-specifying $L$ is a noisy covariance estimate.

\subsubsection{Dimson lead-lag}

The design is augmented with $L$ lagged copies of the factor panel and the
coefficients are summed factor by factor (Section \ref{sec:dimson}). Only lags are
included; a lead would use tomorrow's factor return to explain today's security
return, which is unavailable at forecast time.

The standard error of the sum uses the full coefficient covariance:

\begin{equation}
\operatorname{se}\!\left(\hat\beta^{\mathrm{Dimson}}_j\right)
= \sqrt{\textstyle\sum_{\ell}\sum_{m} \operatorname{Cov}\!\left(\hat\beta^{(\ell)}_j, \hat\beta^{(m)}_j\right)} .
\end{equation}

The lag coefficients are typically negatively correlated, so the diagonal-only
approximation $\sqrt{\sum_\ell \operatorname{se}^2}$ overstates the uncertainty. It is
used only as a conservative fallback when no covariance is supplied, and the caller is
told to treat the resulting interval as wide.

\subsubsection{Collinearity diagnostics}

The \textbf{condition number} is the ratio of the largest to the smallest singular
value of the standardised design matrix. It is the invariant that governs how far
$\beta^{\prime}\Sigma\beta$ can be amplified by near-collinearity, and it is what the
reliability gate of Section \ref{sec:reliability} reads.

\textbf{Variance inflation factors} are computed from the inverse correlation matrix,

\begin{equation}
\mathrm{VIF}_j \;=\; \left[R^{-1}\right]_{jj},
\qquad R = \operatorname{corr}(X),
\end{equation}

which is algebraically identical to $1/(1-R_j^2)$ from $k$ auxiliary regressions but
costs one matrix inversion rather than $k$ least-squares fits --- a material difference
at $k=40$ over a thousand rolling windows. The auxiliary-regression form is retained
in the code as the reference implementation, and a test asserts the two agree.

\subsubsection{Beta stability}

Between consecutive windows, on the factors common to both:

\begin{equation}
d_{L_1} = \sum_{j \in C} \left|\beta_j^{(t)} - \beta_j^{(t-1)}\right|,
\qquad
\rho = \operatorname{corr}\!\left(\beta^{(t)}_C, \beta^{(t-1)}_C\right),
\qquad |C| \text{ recorded.}
\end{equation}

Recording $|C|$ is not optional: $d_{L_1}$ falls mechanically as the common basis
narrows, so the two must be read together.

% ---------------------------------------------------------------------
\subsection{Orthogonalisation}
\label{app:orth}

Given a target series $y$ and a regressor block $X$ drawn from the levels above it in
the hierarchy, the rolling residual is computed as follows.

\begin{enumerate}[leftmargin=1.6em,itemsep=0.3ex]
\item Coefficients are refitted at $t$ whenever $t$ has reached the next refit point,
  on the trailing window $[\max(0,\,t-504),\,t)$, requiring at least
  $\max(252,\ k+2)$ complete observations. The window is half-open: the fit uses data
  strictly before $t$.
\item The next refit point is set to $t + 21$.
\item The residual at $t$ is $y_t - \hat a - \hat b^{\prime}X_t$, using whichever
  coefficients are currently in force.
\item Before the first successful fit, the residual is undefined rather than equal to
  the raw series --- a factor's orthogonalised history starts when it can be
  estimated, not before.
\end{enumerate}

Two variants exist for comparison and are exercised by the test suite: an expanding
window, and an exact full-sample fit. Table \ref{tab:causal} reports what each costs.
The full-sample variant is the only one that achieves exact orthogonality, and it is
the only one that uses future information; it is never used to build the panel.

A Gram-Schmidt routine is also available, which orthogonalises a panel's columns
sequentially left to right. It is used for diagnostics rather than construction,
because the block hierarchy --- not column order --- is what encodes the economic
argument about which factor should absorb shared variation.

% ---------------------------------------------------------------------
\subsection{Covariance estimation}
\label{app:cov}

\subsubsection{Sample}

The unbiased sample covariance on complete rows, annualised by 252. Included as a
baseline and as the reference against which shrinkage is tested; it is not a
sensible choice at $K=39$ on a rolling window.

\subsubsection{Exponentially weighted}

RiskMetrics-style, with weights $w_t \propto 2^{-(T-t)/H}$ normalised to sum to one
and a default half-life of $H=60$ days. Two details matter.

Deviations are taken from the \emph{weighted} mean rather than the simple mean, so
the estimate is internally consistent. And the result is bias-corrected for weighted
sampling,

\begin{equation}
\hat\Sigma = \frac{\sum_t w_t (x_t-\bar x_w)(x_t-\bar x_w)^{\prime}}{1 - \sum_t w_t^2},
\end{equation}

since $1/\sum_t w_t^2$ is the effective sample size when the weights sum to one.

\begin{decision}{Sixty days, not eleven}
The classic RiskMetrics $\lambda = 0.94$ corresponds to a half-life of about 11 days.
That is too jumpy for a risk model used to set limits: the forecast moves materially
on a single day's return. Sixty days is the usual compromise and is the default here.
\end{decision}

\subsubsection{Ledoit-Wolf shrinkage}

Shrinkage toward a \emph{constant-correlation} target \citep{ledoit2003}. The target
$F$ keeps each factor's own variance and replaces every correlation by the average
correlation $\bar r$:

\begin{equation}
F_{ij} = \bar r\, s_i s_j \;\; (i \neq j), \qquad F_{ii} = s_i^2 .
\end{equation}

This is a considerably better description of a factor panel than the identity target
used in the simpler version of the method, which asserts that the factors are
uncorrelated --- precisely the assumption a factor panel violates.

The optimal intensity is

\begin{equation}
\delta^{\star} = \max\!\left\{0,\ \min\!\left\{1,\ \frac{\pi - \rho}{\gamma}\cdot\frac{1}{T}\right\}\right\},
\qquad \hat\Sigma = \delta^{\star} F + (1-\delta^{\star}) S,
\end{equation}

where $\pi$ is the summed asymptotic variance of the sample covariance entries, $\rho$
the asymptotic covariance between the target's estimation error and the sample's, and
$\gamma = \|F - S\|_F^2$ the squared misspecification of the target. The derivation
assumes the maximum-likelihood scaling for $S$ (divisor $T$), which is what the
implementation uses internally; the result is converted to the unbiased scaling
(divisor $T-1$) before being returned, so that it is directly comparable with the
sample estimator. The realised $\delta^{\star}$ is stored on the result --- it is
informative in itself, since a high intensity says the sample estimate was carrying
little signal.

\subsubsection{Blend}

$\Sigma = a\,\Sigma_{\mathrm{EWMA}} + (1-a)\,\Sigma_{\mathrm{shrunk}}$, default
$a = 0.6$. The EWMA leg supplies responsiveness to the current regime; the shrunk leg
supplies conditioning. This is the default estimator.

\subsubsection{Positive semi-definiteness}

Every estimate is checked. Where the smallest eigenvalue is negative, the matrix is
repaired by clipping eigenvalues at a small positive floor and reconstructing, then
\emph{rescaled} so that the diagonal returns to the original factor variances.

The rescale is the part that is easy to omit and important to keep: clipping alone
raises the diagonal, which raises every volatility computed from the matrix. The
repair is recorded on the result rather than applied silently.

\subsubsection{Specific risk}

\begin{equation}
\sigma^2_\varepsilon \;=\; \max\!\left\{\sigma^2_{\mathrm{floor}},\;
(1-k)\,\mathrm{EWMA}_H\!\left[\varepsilon_t^2\right] \cdot 252 \;+\; k\,\sigma^2_{\mathrm{peer}}\right\},
\end{equation}

with $H = 60$ days, $k = 0.25$, and $\sigma_{\mathrm{floor}} = 2\%$ annualised. The
peer variance is the cross-sectional median residual volatility across every
instrument estimated under the same specification, squared. Both the raw and the
floored values are returned, along with a flag recording whether the floor bound,
so that a floored number is never mistaken for a measured one.

\subsubsection{Assembly and attribution}

\begin{equation}
\Sigma_r = B\,\Sigma_F\,B^{\prime} + \operatorname{diag}(\sigma^2_\varepsilon),
\qquad
\sigma_p = \sqrt{\beta^{\prime}\Sigma_F\beta + \sigma^2_\varepsilon} .
\end{equation}

Risk contributions use the standard Euler decomposition,
$\mathrm{RC}_j = \beta_j\,(\Sigma_F\beta)_j / \sigma_p$, which sums to the factor part
of total risk exactly. Contributions are aggregated to block level for display,
because a forty-row contribution table is not a decomposition anyone reads.

% ---------------------------------------------------------------------
\subsection{Forecast evaluation}
\label{app:eval}

\subsubsection{Realised volatility and alignment}

For horizon $h$, realised volatility at $t$ is the annualised standard deviation of
$r_{t+1},\dots,r_{t+h}$. The forecast it is compared against uses loadings from a
window ending at or before $t$ and a covariance estimated on returns up to $t$. The
forward window is constructed by an explicit shift rather than by a rolling operation
with a default alignment, because an off-by-one in this direction is lookahead and
produces a flatteringly good result.

\subsubsection{Bias statistics}

Two are computed. The \textbf{mean bias ratio} is
$\overline{\sigma^{\mathrm{real}}/\sigma^{\mathrm{pred}}}$; it is intuitive but noisy,
being a ratio of volatilities and therefore right-skewed. The \textbf{bias statistic}
is the sharper and is what the text quotes:

\begin{equation}
z_t = \frac{r_t}{\sigma^{\mathrm{pred}}_t},
\qquad
\mathrm{bias} = \operatorname{sd}(z_t),
\end{equation}

the standard deviation of returns divided by the volatility forecast in force at
each date. It equals 1 when the model is calibrated and exceeds 1 when risk is
underestimated. The kurtosis of $z$ is recorded alongside it: a bias statistic of 1
with a kurtosis of 8 describes a model that is right on average and wrong in the
tails.

\subsubsection{Mincer-Zarnowitz}

Realised variance regressed on predicted variance \citep{mincer1969}:

\begin{equation}
\left(\sigma^{\mathrm{real}}_t\right)^2 = a + b\left(\sigma^{\mathrm{pred}}_t\right)^2 + u_t,
\qquad H_0:\ (a,b) = (0,1).
\end{equation}

Run on variances rather than volatilities, because variance is what the model
forecasts and what aggregates linearly. Individual $t$-tests for $a=0$ and $b=1$ are
reported together with a joint Wald test,
$W = (\hat\theta - \theta_0)^{\prime} V^{-1} (\hat\theta - \theta_0) \sim \chi^2_2$.

HAC standard errors are not optional here. Overlapping realised windows --- a 21-day
forward window observed daily --- induce strong serial correlation in $u_t$ by
construction, and ordinary standard errors would reject far too often. Section
\ref{sec:risk} gives the separate caveat on interpreting the slope.

\subsubsection{Value at risk and expected shortfall}

\begin{equation}
\mathrm{VaR}_t(\alpha) = -z_\alpha\,\frac{\sigma^{\mathrm{pred}}_t}{\sqrt{252}},
\end{equation}

under either a Gaussian quantile or a scaled Student-$t$ quantile with 5 degrees of
freedom, the latter being the more honest default for daily returns. Expected
shortfall is computed at the 97.5\% level under the same distributional choice.

\subsubsection{Kupiec proportion of failures}

\citep{kupiec1995} A likelihood ratio on the exception \emph{count}, $\chi^2_1$:

\begin{equation}
\mathrm{LR}_{\mathrm{POF}} = -2\left[(T-x)\log(1-p) + x\log p - (T-x)\log(1-\hat\pi) - x\log\hat\pi\right],
\end{equation}

with $p = 1-\alpha$, $x$ exceptions in $T$ observations and $\hat\pi = x/T$. At the
boundaries $x = 0$ or $x = T$ the ratio is degenerate; the implementation falls back
to an exact two-sided binomial probability rather than returning a spurious statistic.

\subsubsection{Christoffersen conditional coverage}

\citep{christoffersen1998} $H_0$: the correct rate \emph{and} independence.
Transitions are counted as $n_{ij}$ (from state $i$ to state $j$, breach = 1), and

\begin{equation}
\mathrm{LR}_{\mathrm{CC}} = \mathrm{LR}_{\mathrm{POF}} + \mathrm{LR}_{\mathrm{IND}} \sim \chi^2_2 .
\end{equation}

Where the transition counts cannot identify the independence term --- no breach ever
follows a non-breach, or one of the states is never visited --- the implementation
reports Kupiec alone on one degree of freedom rather than pretending to test
something the data cannot support.

This is the test that matters most and is reported even when it fails, because
clustered breaches are the dangerous failure: a model can produce exactly the right
number of exceptions over a decade and place them all in one week.

% ---------------------------------------------------------------------
\subsection{Density estimation}
\label{app:density}

The distribution panel exists to answer one question --- how far from Gaussian is this
series, and where --- so its defaults are chosen to keep the tails visible rather than
to produce a smooth picture.

\paragraph{Robust scale.} The smaller of the sample standard deviation and
$\mathrm{IQR}/1.349$. The two coincide for a Gaussian sample; on a fat-tailed series
the second is much smaller, and taking the minimum ties bandwidth and bin width to the
bulk of the distribution, which is where resolution is wanted.

\paragraph{Bin count.} Freedman-Diaconis \citep{freedman1981},
$h = 2\,\mathrm{IQR}\,n^{-1/3}$, with the resulting count clamped to $[20, 220]$ for
legibility. Scott's rule uses the standard deviation and would be inflated by exactly
the outliers that create the problem.

\paragraph{Bandwidth.} Silverman's rule \citep{silverman1986} on the robust scale,
$h = 0.9\,s\,n^{-1/5}$. Using the plain standard deviation here over-smooths a
fat-tailed series into a flat blob.

\paragraph{Overlays and range.} A Gaussian fit and a Student-$t$ fit are drawn over
the empirical density; the comparison between them is the whole purpose of the panel.
The drawing range is the 0.1st to 99.9th percentile, and the observations outside it
are \emph{counted and reported}, not silently dropped --- a plot that hid its own tail
would be worse than the naive one it replaces.

\newpage
% =====================================================================
\section{Factor registry}
\label{app:registry}

The forty factors, their block, their hierarchy level and the factors each is
residualised against. Level 0 factors are used raw. A factor's orthogonalisation
targets are always factors at a lower level, so the hierarchy is acyclic by
construction and the panel can be built in one pass.

The construction of each factor --- the exact formula over named instruments --- is
given in the companion document \emph{factor-formulas}, which is generated from this
same registry.

% Generated by scripts/build_paper.py — do not edit by hand.
\begingroup\small
\setlength{\LTleft}{0pt}\setlength{\LTright}{\fill}
\begin{longtable}{@{}llc>{\raggedright\arraybackslash}p{0.40\textwidth}@{}}
\toprule
\textbf{Factor} & \textbf{Block} & \textbf{Level} & \textbf{Residualised against} \\
\midrule
\endfirsthead
\toprule
\textbf{Factor} & \textbf{Block} & \textbf{Level} & \textbf{Residualised against} \\
\midrule
\endhead
\bottomrule
\endfoot
\texttt{\footnotesize eq\_global} & Equity Market & 0 & \textit{--- used raw} \\
\texttt{\footnotesize eq\_cyclical} & Equity Market & 1 & \texttt{\footnotesize eq\_global} \\
\texttt{\footnotesize eq\_em} & Equity Market & 1 & \texttt{\footnotesize eq\_global} \\
\texttt{\footnotesize eq\_europe} & Equity Market & 1 & \texttt{\footnotesize eq\_global} \\
\texttt{\footnotesize eq\_japan} & Equity Market & 1 & \texttt{\footnotesize eq\_global} \\
\texttt{\footnotesize eq\_size} & Equity Market & 1 & \texttt{\footnotesize eq\_global} \\
\texttt{\footnotesize eq\_us} & Equity Market & 1 & \texttt{\footnotesize eq\_global} \\
\addlinespace
\texttt{\footnotesize sty\_lowvol} & Equity Style & 2 & \texttt{\footnotesize eq\_global}, \texttt{\footnotesize eq\_us}, \texttt{\footnotesize eq\_cyclical} \\
\texttt{\footnotesize sty\_momentum} & Equity Style & 2 & \texttt{\footnotesize eq\_global}, \texttt{\footnotesize eq\_us}, \texttt{\footnotesize eq\_cyclical} \\
\texttt{\footnotesize sty\_quality} & Equity Style & 2 & \texttt{\footnotesize eq\_global}, \texttt{\footnotesize eq\_us}, \texttt{\footnotesize eq\_cyclical} \\
\texttt{\footnotesize sty\_value} & Equity Style & 2 & \texttt{\footnotesize eq\_global}, \texttt{\footnotesize eq\_us}, \texttt{\footnotesize eq\_cyclical} \\
\addlinespace
\texttt{\footnotesize rt\_breakeven} & Rates & 0 & \texttt{\footnotesize rt\_us\_level} \\
\texttt{\footnotesize rt\_ea\_level} & Rates & 0 & \texttt{\footnotesize rt\_us\_level} \\
\texttt{\footnotesize rt\_jp\_level} & Rates & 0 & \texttt{\footnotesize rt\_us\_level} \\
\texttt{\footnotesize rt\_uk} & Rates & 0 & \texttt{\footnotesize rt\_us\_level} \\
\texttt{\footnotesize rt\_us\_curve} & Rates & 0 & \textit{--- used raw} \\
\texttt{\footnotesize rt\_us\_level} & Rates & 0 & \textit{--- used raw} \\
\texttt{\footnotesize rt\_us\_slope} & Rates & 0 & \textit{--- used raw} \\
\addlinespace
\texttt{\footnotesize cr\_em\_hard} & Credit & 1 & \texttt{\footnotesize rt\_us\_level}, \texttt{\footnotesize rt\_us\_slope}, \texttt{\footnotesize cr\_hy}, \texttt{\footnotesize eq\_global} \\
\texttt{\footnotesize cr\_hy} & Credit & 1 & \texttt{\footnotesize rt\_us\_level}, \texttt{\footnotesize rt\_us\_slope}, \texttt{\footnotesize eq\_global} \\
\texttt{\footnotesize cr\_ig} & Credit & 1 & \texttt{\footnotesize rt\_us\_level}, \texttt{\footnotesize rt\_us\_slope} \\
\texttt{\footnotesize cr\_loans} & Credit & 1 & \texttt{\footnotesize rt\_us\_level}, \texttt{\footnotesize cr\_hy} \\
\texttt{\footnotesize cr\_em\_local} & Credit & 2 & \texttt{\footnotesize rt\_us\_level}, \texttt{\footnotesize cr\_em\_hard}, \texttt{\footnotesize fx\_usd} \\
\addlinespace
\texttt{\footnotesize fx\_usd} & FX & 0 & \textit{--- used raw} \\
\texttt{\footnotesize fx\_carry} & FX & 1 & \texttt{\footnotesize fx\_usd}, \texttt{\footnotesize fx\_jpy} \\
\texttt{\footnotesize fx\_em} & FX & 1 & \texttt{\footnotesize fx\_usd} \\
\texttt{\footnotesize fx\_jpy} & FX & 1 & \texttt{\footnotesize fx\_usd} \\
\addlinespace
\texttt{\footnotesize cm\_broad} & Commodities & 0 & \textit{--- used raw} \\
\texttt{\footnotesize cm\_agri} & Commodities & 1 & \texttt{\footnotesize cm\_broad} \\
\texttt{\footnotesize cm\_energy} & Commodities & 1 & \texttt{\footnotesize cm\_broad} \\
\texttt{\footnotesize cm\_gold} & Commodities & 1 & \texttt{\footnotesize cm\_broad} \\
\texttt{\footnotesize cm\_industrial} & Commodities & 1 & \texttt{\footnotesize cm\_broad} \\
\addlinespace
\texttt{\footnotesize vol\_equity} & Volatility & 1 & \texttt{\footnotesize eq\_global} \\
\texttt{\footnotesize vol\_rates} & Volatility & 1 & \texttt{\footnotesize rt\_us\_level} \\
\texttt{\footnotesize vol\_variance\_premium} & Volatility & 1 & \texttt{\footnotesize eq\_global} \\
\addlinespace
\texttt{\footnotesize liq\_conditions} & Liquidity \& Stress & 0 & \textit{--- used raw} \\
\texttt{\footnotesize liq\_funding} & Liquidity \& Stress & 0 & \textit{--- used raw} \\
\texttt{\footnotesize liq\_risk\_off} & Liquidity \& Stress & 2 & \texttt{\footnotesize eq\_global}, \texttt{\footnotesize rt\_us\_level}, \texttt{\footnotesize cr\_hy} \\
\addlinespace
\texttt{\footnotesize arp\_trend} & Alternative Risk Premia & 3 & \texttt{\footnotesize eq\_global}, \texttt{\footnotesize rt\_us\_level}, \texttt{\footnotesize cm\_broad}, \texttt{\footnotesize fx\_usd} \\
\texttt{\footnotesize arp\_xs\_momentum} & Alternative Risk Premia & 3 & \texttt{\footnotesize eq\_global}, \texttt{\footnotesize arp\_trend} \\
\end{longtable}
\endgroup

